\cotefile{162-5}
\scanpage{20}
\begin{center}
PUBLICATIONS OF DANIEL QUILLEN
\end{center}

1.\quad Spectral sequences of a double semi-simplicial group. Topology 5 (1966) 155-157.

2.\quad \uline{Homotopical} \uline{Algebra}. Lecture Notes in Mathematics No. 43. Springer-Verlag, Berlin-N.Y. 1967.

3.\quad The geometric realization of a Kan fibration is a Serre fibration. Proc. Amer. Math. Soc. 19 (1968) 1499-1500.

4.\quad On the representation of hermitian forms as sums of squares. Invent. Math. 5 (1968), 237-242.

5.\quad On the associated graded ring of a group ring. J. Algebra 10 (1968) 411-418.

6.\quad Some remarks on etale Homotopy theory and a conjecture of Adams, Topology 7 (1968), 111-116.

7.\quad Rational Homotopy Theory, Annals of Math. v. 90 (1969).

8.\quad The Adams Conjecture, Preprint (1970).

9.\quad An application of simplicial profinite groups.Comment. Math. Helv. 44 (1969) 45-60.

10.\quad On the endomorphism ring of a simple module over an enveloping algebra. Proc. Amer. Math. Soc. 21 (1969)171-172.

11.\quad The sprectrum of an equivariant cohomology ring. (To appear).

% archive sheet number: 19

\scanpage{22}
\begin{flushright}
\uline{tapis Quillen} (Notes du 6.9.68)
\end{flushright}

1.\quad Catégories formelles.

Catégorie formelle \inserted{$\mathcal C$} ``sur $X_0$'' i.e. dont \unsure{le 'objet|l'objet} des objets est le schéma $X_0$ (regardé comme \deleted{\unsure{ind-}} schéma formel i.e. \uline{ind-schéma} particulier).

Revient à la donnée d'un pro-Anneau strict \inserted{``admissible''} $A$ sur $X_0$, augmenté vers $\underline O_{X_0}${\reversemarginpar\margin{%
% Diagram in the left margin beside « augmenté vers ».
\diagram{\begin{tikzcd}[ampersand replacement=\&,row sep=large]
X_2 \arrow[d,shift right=4.4ex] \arrow[d,shift right=2.2ex] \arrow[d] \\
X_1 \arrow[u,shift right=2.3ex,bend right=9] \arrow[u,shift right=4.4ex,bend right=9] \arrow[d,shift right=2ex] \arrow[d] \\
X_0 \arrow[u,shift right=2.4ex,bend right=10]
\end{tikzcd}}
}}, l'idéal d'augmentation étant idéal de définition, avec deux structures de $\underline O_{X_0}$-modules (g. et dr.), et un hom. \inserted{de bi-$\underline O_{X_0}$-proanneaux} $A\xrightarrow{\Delta}A_{\mathrm{dr}}\otimes_{\underline O_{X_0}}{}_{\mathrm g}A\ \deleted{\unsure{\longrightarrow}}$, satisfaisant des axiomes évidents a) ``associativité'' de $\Delta$ b) unitarité g. et dr. vis à vis de l'augmentation. [Cela sera-t-il nécessairement \unsure{un|une} \deleted{catégorie} \uline{groupoïde} formel, i.e. y-aura-t-il une \uline{symétrie} $s$ de $A$, jouant vis-à-vis de $\Delta$ et de l'augmentation $\varepsilon$ le rôle \unsure{du|de} passage à l'inverse ?]

On pose $\Omega_{\mathcal C}={}$faisceau conormal de $X_0$ dans $X_1$ ; à priori c'est un pro-faisceau qu. ch. sur $X_0$, mais dans le cas où $A$ est $I$-\uline{adique} [$I$ l'idéal d'augmentation] c'est un module qu. ch. ordinaire.

On définit alors sur $\bigwedge^{*}\Omega$ une différentielle d'anneau gradué \inserted{$d:d_{\mathcal C}$} Sa valeur sur $\bigwedge^0\Omega=\underline O_{X_0}$ est
\[
df=\mathrm{pr}_2^{*}(f)-\mathrm{pr}_1^{\unsure{\alpha|*}}(f)\mod I^2
\]

% slanted
\margin{\uline{\unsure{n'interviendra|n'intervient|infinitésimale}} \uline{\unsure{pas|voir|vu} plus bas}}
\uline{Suite à développer} : $A=(A_\alpha)$, regardons $A_\alpha$ comme module pour la structure gauche, alors l'application $f\longmapsto\mathrm{pr}_2^{*}(f):\underline O_X\mathbin{-} A_\alpha$ est un op. diff. [d'ordre $n$ si $I_\alpha^{n+1}=0$], et par suite tte \deleted{\unsure{forme li}} section de $\underline{\mathrm{Hom}}_{\underline O_{\unsure{\alpha|X}}}(A_\alpha,\underline O_X)=\underline{\mathscr D}_{\alpha}^{\mathcal C}$ définit un op. différentiel dans $\underline O_X$. Alors $\mathscr D_{\unsure{0|\mathcal C}}^{\mathcal C}=\underline{\lim}_{\alpha}\mathscr D_{\alpha}^{\mathcal C}$ est un faisceau d'algèbres \inserted{\unsure{unitaires n.c.|unitaires ass.}} [grâce à l'application $\Delta$], qui s'envoie dans le faisceau des op différentiels [qui correspond

% archive sheet number: 21

\scanpage{23}
au \unsure{fait|cas|fond} où $\mathcal C$ serait la formalisation de la catégorie triviale définie par $X_0/\mathrm{Spec}\,Z$). En particulier
\[
\mathscr T^{\mathcal C}=\deleted{\unsure{\mathrm{dual}|\text{déf}}}\ \underline{\mathrm{Hom}}_{\underline O_{X_0}}(\Omega_{\mathcal C}^{\deleted{\unsure{\mathcal C}}},\underline O_X)=\varinjlim_{\alpha}\underline{\mathrm{Hom}}_{\underline O_{X_0}}(\Omega_{\unsure{G_\alpha|\mathcal C_\alpha}}^{\mathcal C},\underline O_{X_0})
\]
opère sur $\underline O_{X_0}$ par \uline{dérivations}\,\ldots. Cf aussi plus bas pour sa déduction à partir du complexe de De Rham $\Omega_{\mathcal C}^{*\deleted{\unsure{\mathcal C}}}$ de $\mathcal C$.

2.\quad \uline{Complexes de De Rham} \uline{généraux}.

Appelons complexe de De Rham construit sur un \deleted{\unsure{sché}} Module \inserted{qu. ch.} $\Omega$ sur un schéma $X_0$ [NB marchera aussi pour tout espace annelé, de même que la suite du 1.] \inserted{tout complexe obtenu par} une \deleted{\unsure{op}} différentielle d'anneau gradué sur $\bigwedge^{*}\Omega$.

\uline{Remarque} 1\quad Cette différentielle est connue quand on la connaît en degrés 0 et 1
\[
\begin{aligned}
\delta_0:\underline O_{X_0}\longrightarrow\Omega\quad &\inserted{\alpha)\ \text{\uline{dérivation}}}\quad [\,\text{i.e.}\quad\Omega^1_{X_0/\mathbb Z}\xrightarrow{u}\Omega\\
\inserted{\alpha)\ \delta_0(fg)=\delta_0(f)g+f\delta_0(g)}\quad &\text{par}\quad\delta(f)=u(df)\quad ]\\
\delta_1:\Omega\longrightarrow\bigwedge^2\Omega &
\end{aligned}
\]
Condition nécessaire sur $\delta_1$
\[
\left\{\begin{aligned}
\beta)\ &\delta_1(f\omega)=\delta_0(f)\,\delta_1(\omega)+f\,\delta_1(\omega)\\
\gamma)\ &\delta_1\delta_0=0
\end{aligned}\right.
\]

De plus, on peut choisir $\delta_0$ et $\delta_1$ arbitrairement, de façon à satisfaire ces conditions : ils se prolongent (de façon

% archive sheet number: 22

\scanpage{24}
unique) en une \deleted{\unsure{homo}}\inserted{\uline{anti}}dérivation d'Algèbre\inserted{grâce \deleted{à la première} \inserted{aux} conditions \deleted{\unsure{$\beta)$}}\inserted{$\alpha,\beta$}}, et cette dérivation est de carré nul grâce à la condition $\gamma$.

% slanted
\margin{dégager un lemme \unsure{circonstancié}\,\ldots.}

3.\quad \deleted{Algèbre de} \uline{Anneau de Lie} \uline{associé à un complexe de De Rham}

Soit $\mathscr T=\underline{\mathrm{Hom}}(\Omega,\underline O_{X_0})$, c'est un ind-module quasi-ch. sur $X_0$, qu'on identifie à un Module qu. coh.

Ceci définit, si $X,Y$ sont des sections de $\mathscr T$ \deleted{\unsure{sur}} sur un $U$
\[
[X,Y]\deleted{\unsure{(\omega)|(U)}}\in\Gamma(U,\mathscr T)\qquad\text{avec}
\]
\[
\left\langle\underset{\substack{\rotatebox{-90}{$\in$}\\\Gamma(U,\Omega)}}{\omega},[X,Y]\right\rangle\deleted{\unsure{(\omega)}}=-\langle\delta\omega,X\wedge Y\rangle\deleted{\unsure{-}}+\Theta_X\langle\omega,Y\rangle\deleted{\unsure{+}}-\Theta_Y\langle\omega,X\rangle
\tag*{$(x)$}
\]
où l'action $X\mapsto\Theta_X$ de $\mathscr T$ sur $\underline O_{X_0}$ \inserted{(par dérivations)} est définie par
\[
\begin{aligned}
\Theta_X(f)&=\langle\delta f,X\rangle,\qquad [\,\text{NB}\quad\boxed{\Theta_{gX}=g\,\Theta_X}\,]\ \text{résulte de rien}\\
&\hspace{6em}(\text{résulte de }\alpha))\ \boxed{\Theta_X(fg)=\Theta_X(f)g+f\Theta_X(g)}
\end{aligned}
\]
On constate \inserted{grâce à $\alpha)$ et $\beta)$} que le deuxième mb de $(x)$ est bien \uline{$\underline O_{X_0}$-linéaire} en $\omega$, donc définit une section $X\wedge Y$.

De plus, on aura, utilisant $\gamma)$
\[
\deleted{\unsure{[X,Y]}}\ \boxed{\Theta_{[X,Y]}=[\Theta_X,\Theta_Y]}\overset{\mathrm{\unsure{dfn|defn|df}}}{=}\Theta_X\Theta_Y\deleted{\unsure{f}}-\Theta_Y\Theta_X
\tag*{$(xx)$}
\]
\deleted{[\unsure{Si on|Sinon} \unsure{veut|vient} l'autre définition}

[Sans $\gamma$, la \deleted{\unsure{valeur des}} différence des deux mbs sur $f$ est $-\langle\delta\delta f,X\wedge Y\rangle$ ; donc si $\Omega$ loc libre, la condition $\delta_1\delta_0=0$ est \uline{nécessaire} pour qu'on ait $(xx)$]. Ceci vérifie aussi la loi de Lie
\margin{\textcircled{\scriptsize trivial}}
\[
\boxed{[X,X]=0},\qquad\boxed{[[X,Y],Z]+[[Y,Z],X]+[[Z,X],Y]=0}
\]

% archive sheet number: 23

\scanpage{25}
\[
\boxed{[X,\unsure{d|f}Y]=\Theta_X(f)Y+f[X,Y]}\qquad\text{résulte de }\alpha
\]

Prouvons enfin la formule du crochet. Soit $\omega\in\unsure{\Gamma\Omega|\Omega|\Gamma(\Omega)}$, on doit \deleted{vérifier} calculer
\[
\begin{aligned}
&\sum_{\substack{\text{perm.}\\\text{circ. sur}\\X,Y,Z}}
\bigl[-\langle\delta\omega,[X,Y]\wedge Z\rangle+\Theta_{[X,Y]}\langle\omega,\deleted{\unsure{Y}}Z\deleted{\unsure{\omega}}\rangle-\Theta_Z\deleted{\unsure{i_{[X,Y]}}}\langle\omega,[X\unsure{,}\,Y]\rangle\bigr]\\
={}&\sum\bigl[-\langle i_Z\delta\omega,[X,Y]\rangle+(\Theta_X\Theta_Y-\Theta_Y\Theta_X)\langle\omega,Z\rangle+\Theta_Z\langle\delta\omega,X\wedge Y\rangle\\
&\hspace{14em}-\Theta_Z\Theta_X\langle\omega,Y\rangle+\Theta_Z\Theta_Y\langle\omega,X\rangle\bigr]\\
={}&\sum-\langle i_Z\delta\omega,[X,Y]\rangle+\Theta_Z\langle\delta\omega,X\wedge Y\rangle\\
={}&\sum\langle\delta\deleted{\unsure{\Theta_Z}}i_Z\delta\omega,X\wedge Y\rangle\,
\deleted{\unsure{+}}-\underset{\smash{\rule{.35pt}{2pt}\rule{9em}{.35pt}\rule{.35pt}{2pt}}}{\Theta_X i_Y\deleted{\unsure{\delta}}i_Z\delta\omega+\Theta_Y i_X i_Z}\delta\omega+\Theta_Z\langle\delta\omega,X\wedge Y\rangle\\
={}&\sum\bigl(i_X i_Y\deleted{\unsure{\delta i_Z}}\inserted{\delta i_Z}+\Theta_Z\deleted{\unsure{\delta}}i_X i_Y\deleted{\unsure{\delta}}\deleted{\unsure{+}}-\Theta_X i_Y i_Z+\Theta_Y i_X i_Z\bigr).(\delta\omega)
\end{aligned}
\]

% The crossed panel extends from here to the curtailed = −9… + line.
\deleted{On va utiliser} $\deleted{\boxed{\Theta_{\unsure{Z|X}}=i_X\delta+\delta i_X}}$ \deleted{et les formules analogues en $Y,Z$, \unsure{pour trouver|qui donnent}} \deleted{\unsure{permu|poser}}
\[
\begin{aligned}
&\deleted{\sum i_Y i_X\delta i_Z+i_Z\underset{\smash{\rule{.35pt}{2pt}\rule{1.4em}{.35pt}\rule{.35pt}{2pt}}}{\delta i_Y}i_X+\unsure{\delta i_X i_Y i_Z}}\
\deleted{\inserted{\delta i_Z i_Y i_X}-i_X\underset{\smash{\rule{.35pt}{2pt}\rule{1.4em}{.35pt}\rule{.35pt}{2pt}}}{\delta i_Y}i_Z-\delta i_X i_Y i_Z}\\
&\hspace{12em}\deleted{\unsure{+}+i_Y\delta\underset{\smash{\rule{.35pt}{2pt}\rule{1.4em}{.35pt}\rule{.35pt}{2pt}}}{i_X i_Z}+\delta i_Y i_X i_Z}\\
&\deleted{\sum i_Y i_X\delta i_Z+(i_Z\delta i_Y i_X+i_Y\delta i_X i_Z\unsure{+}-i_X\delta i_Y i_Z)-3\delta i_X i_Y i_Z}\\
&\deleted{=-9\,\delta i_X i_Y i_Z+}
\end{aligned}
\]

On constate \unsure{trivialement|bêtement|similairement}
\[
\begin{aligned}
&\sum-\Theta_X i_Y i_Z+\Theta_Y i_X i_Z=-2\sum\Theta_{\unsure{Z|X}}i_X i_Y\\
&\deleted{\sum\Theta_Z i_X i_Y=\unsure{+|-}\sum\Theta_X i_Y i_Z}
\end{aligned}
\]
donc on trouve
\[
\begin{aligned}
(\text{Jacobi})(\omega)&=T(\delta\omega),\quad\text{avec}\\
T&=\sum i_X i_Y\delta i_Z-\deleted{\unsure{\delta}}\overset{\inserted{(\Theta_Z=i_Z\delta}}{\Theta_Z}\,i_X i_Y\ \deleted{\unsure{=\sum i_X i_Y}}
\end{aligned}
\]

% archive sheet number: 24

\scanpage{26}
Utilisons
\[
\boxed{\Theta_X=i_X\delta+\delta i_X}\quad\inserted{\text{(en degré quelconque)}}\quad\text{cf plus bas}
\]
il vient alors, au choix
\[
\begin{aligned}
T&=\sum i_X i_Y\delta i_Z-i_Z\delta i_X i_Y\deleted{\unsure{-\delta i_Z i_X i_Y}}\qquad \unsure{!|/}\quad\text{ou}\\
T(\delta\omega)&=\sum[i_{X\wedge Y},\Theta_Z](\delta\omega)
\end{aligned}
\]
\margin{pourquoi cela est-il nul ??}

On a donc une \uuline{$\Theta$-Algèbre de Lie} sur $X_0$, i.e. un \uline{Anneau de Lie} \deleted{sur} $\mathscr T$ sur $\deleted{\unsure{\underline O_{X_0}}}X_0$, muni d'un \uline{hom} \inserted{$\Theta:X\longmapsto\Theta_X$} d'\uline{Anneaux} de Lie dans $\mathscr T_{X_0/Z}$, et d'une structure de $\underline O_{X_0}$-Module rendant $\Theta$ $\underline O_{X_0}$-linéaire et satisfaisant la formule $[\deleted{\unsure{\Theta}}X,fY]=\deleted{\unsure{f\Theta}}f[X,Y]+\Theta_X(f)Y$.

4.\quad \uline{Complexe de Cochaines associé à une $\Theta$-algèbre de Lie} \inserted{$\mathscr T$}

Cf. mes notes Bourbaki ; on pose \deleted{(\unsure{supposant} qu'on soit sur un pt)}
\[
\underline{\unsure{C|\mathcal C}}^p=\underline{\mathrm{Hom}}(\bigwedge^p\mathscr T,O_{X_0})
\]
et on y définit une différentielle par une formule connue\,\ldots\ldots.

Lorsque $\mathscr T$ provient d'un complexe de De Rham $(\bigwedge^*\Omega,\delta)$, alors l'hom d'alg. graduées évident $\bigwedge^*\Omega\longrightarrow C^{\unsure{*|\bullet}}(\mathscr T)$ est \uline{compatible} avec les \deleted{structures} différentielles.

% archive sheet number: 25

\scanpage{27}
5.\quad Les $\Theta$-Algèbres loc libres \inserted{de t.f. $\mathscr T$} sur $X_0$ correspondent exactement aux complexes de De-Rham sur $X_0$ construits sur des $\Omega$ loc. libres de t. f, par les correspondances de $3^\circ$ et $4^\circ$.

6.\quad Revenons au cas d'une catégorie formelle sur $X_0$, définie par $A=(A_\alpha)$ sur $X_0\deleted{\unsure{\mathcal C}}$\,\ldots, et le complexe de De Rham associé $(\Omega^*\unsure{\underset{\nearrow}{=}|=}\bigwedge^*\Omega,\delta)$.

\uline{Th. de Quillen} \inserted{Supposons $X_0$ de car. nulle} Le foncteur $A\longmapsto(\Omega^*,\delta)$ induit une \uline{équivalence} entre la catégorie des catégories formelles \inserted{adiques} sur $X_0$ dont le $\Omega$ \unsure{est|$\to$} \uline{plat}, et la catégorie des \deleted{algèbres}\inserted{complexes} de De Rham $(\bigwedge^{\unsure{\alpha|*}}\Omega,\delta)$ sur $X_0$ dont le $\Omega$ \unsure{soit|est} \uuline{plat}. [Les $A$ en question ont la vertu que $\mathrm{Sym}\,\Omega\longrightarrow\mathrm{Gr}(A)$ est un isomorphisme, ce qui \unsure{s'apparente|signifie} à une propriété de ``lissité différentielle''.

Ceci contient le th. de Cartier sur les groupes formels en car. nulle.

Il y a une variante ``à puiss. divisées'', quand on met sur l'idéal d'augmentation de $A$ \inserted{$=(A_n)_{n\in\mathbb N}$} une structure d'\deleted{\unsure{Algèbre|algèbre}} puiss divisées, qui passe au quotient sur les $A_n$, chaque $I_n$ étant nilpotent d'ordre $n$, et la filtration étant celle provenant des puiss. divisées. [À voir si dans cet énoncé il faut \unsure{imposer|toujours} des relations entre la structure \deleted{\unsure{de}} puiss. divisées, et l'application $\Delta$ \ldots]

% archive sheet number: 26

\scanpage{28}
7.\quad \uline{Coefficients} \inserted{\uline{sur catégorie formelle}} Soit $M$ un $\underline O_{X_0}$-module. La notion de ``donnée de descente'' sur $M$ relativement à $\mathcal C$ [ou ``$\mathcal C$-stratification'', ou ``opération de $\mathcal C$ sur $M$''] est claire. On constate alors que toute telle structure permet de mettre sur
\[
\bigwedge^*\Omega\otimes_{O_{X_0}}M
\]
un opérateur différentiel $\delta_M$, tel que
\[
\delta_M(fx)=\delta(f)x\deleted{\unsure{\delta_M(x)}}+(-1)^{\deg\unsure{f|F}}f\,\delta_M(x)
\tag*{$(*)$}
\]
[Il suffit $=$ d'avoir une donnée de recollement d'ordre 1]

Un tel opérateur est connu quand on connaît
\[
M\xrightarrow{\delta_M^1}\Omega\otimes M
\]
satisfaisant $(*)$ \deleted{\unsure{$\delta_M^1($}} pour $f\in\Gamma\underline O_{X_0}$ (donc $\deg f=0$).

Le fait que l'opérateur déduit \inserted{sur $\bigwedge^*\Omega\otimes M$ (cf plus bas)} soit de carré nul provient de la condition de descente. [il suffit qu'elle soit vraie à l'ordre 2 \ldots??].

8.\quad \uline{Coefficients relatifs à une Algèbre de De Rham}

Plus gén, si $(\bigwedge^*\Omega,\delta)$ est une alg de De Rham sur $X_0$, on \deleted{\unsure{peut}} regarde les modules $M$ avec une \uline{connexion} relativement à celle-ci, i.e. un morphisme
\[
M\overset{\delta_M}{\mathbin{-}}\ \Omega\otimes M
\]
satisfaisant
\[
\delta_M(fx)=\delta(f)\otimes x+f\,\delta_M(x).
\]

% archive sheet number: 27

\scanpage{29}
On peut alors de façon unique prolonger en une \deleted{dérivation} opération $\delta_M$ sur $\bigwedge^*\Omega\otimes M$, en posant
\[
\delta_M(\omega x)=\delta(\omega)\deleted{\unsure{\otimes}}x+\omega\,(\delta_M x)
\qquad\left(\begin{aligned}
&\omega\in\Gamma\bigwedge^*\Omega,\\
&x\in\Gamma\bigwedge^*\Omega\otimes M
\end{aligned}\right)
\]

Pour qu'on ait $\delta_M\delta_M=0$, il f. et s qu'on l'ait \deleted{\unsure{$\delta_M$}} en degré 0, i.e. que le composé
\[
M\xrightarrow{\delta_M^{\unsure{0|\bullet}}}\Omega\otimes M
\overset{\delta_M^1}{\mathbin{\unsure{-|\to}}}\bigwedge^2\Omega\otimes\deleted{\unsure{\wedge}} M
\]
soit nul. Ce composé s'appelle le \uline{tenseur} \uline{courbure}.

La donnée de \inserted{la connexion} $\delta_M^0$ permet de faire opérer $\mathscr T=\underline{\mathrm{Hom}}_{\underline O_X}(\Omega,\underline O_X)$ dans $M$, \unsure{grâce|par|grâce à}.
\[
\Theta_X^M(x)=(X\otimes\mathrm{id}_M)(\delta_M^0(x))
\]
On a alors
\[
\begin{array}{r@{\ }l}
\begin{aligned}a)\\[1ex] b)\end{aligned}&
\left\{\begin{aligned}
\Theta_X^M(fx)&=\Theta_X(f)x+f\Theta_X^M(x)\\[1ex]
\Theta_{fX}^M(x)&=f\Theta_X^M(x)
\end{aligned}\right.
\end{array}
\]
D'autre part, si la connexion est sans \deleted{\unsure{torsion|intégrable}}\inserted{courbure}, alors il est vrai qu'on a
\[
c)\quad\left\{\Theta^M_{[X,Y]}=[\Theta_X^{\unsure{H|M}},\Theta_Y^{H}].\right.
\]

% archive sheet number: 28

\scanpage{30}
9.\quad \deleted{\unsure{Partant} d'une} \uline{Opération d'une $\Theta$-Alg. de Lie}

Soit $\mathscr T$ une $\Theta$-Algèbre de Lie sur $X_0$, $M$ un $\underline O_{X_0}$-Module. On dit que \deleted{$\mathscr T$} $M$ a une $\mathscr T$-connexion si on fait ``opérer'' $\mathscr T$ sur $M$ de façon à satisfaire à a) et b) plus haut.

Ceci définit alors un tenseur courbure
\[
K\in C^2(\mathscr T,\unsure{\mathrm{End}|\mathscr E\mathrm{nd}}_{\underline O_{X_0}}(M))
\]
en posant
\[
K(X,Y)=\Theta_{[X,Y]}^M-[\Theta_X^M,\Theta_Y^{\unsure{H|M}}]
\]
Le tenseur courbure est donc nul ssi on a
\[
\Theta_{[X,Y]}^M=[\Theta_X^M,\Theta_Y^H],
\]
alors la connexion est \uline{dite intégrable}. Si $\mathscr T$ provient de $\Omega$ \inserted{$(\mathscr T=\underline{\mathrm{Hom}}_{\underline O_{X_0}}(\Omega,\underline O_{X_0}))$}, et $M$ \deleted{\unsure{provient}} muni d'une $\Omega$-connexion $\delta_M^0$, d'où un tenseur courbure
\diagram{\begin{tikzcd}[ampersand replacement=\&,column sep=large]
M\arrow[r]\arrow[rr,bend right=22,no head,"K'"'] \& \bigwedge^2\Omega\otimes M\arrow[r] \& C^2(\mathscr T,M)
\end{tikzcd}}
\unsure{si conventions|j'écris maintenant|j'aimerais}
% slanted
\margin{Attention signes !}
\[
K(X,Y)(u)=K'(\unsure{x|u})(X,X).
\]

\uline{10}.\quad Si \inserted{$(\bigwedge^*\Omega,\delta)$ \deleted{\unsure{avec}} défini} $\Omega$ loc. libre de t.f. ($\Longleftrightarrow\mathscr T$ \inserted{$\Theta$-alg. de Lie \unsure{sur|avec} $\mathscr T$} loc. libre de t.f.) alors pour un Module $M$ sur $\underline O_{X_0}$,

% archive sheet number: 29

\scanpage{31}
les deux notions de connexion se correspondent \unsure{biunivoquement}, de $=$ que les deux notions de \inserted{tenseur courbure, de} connexion intégrable \ldots, et bien sûr les complexes
\[
\bigwedge^*\Omega\otimes M\qquad\text{et}\qquad C^*(\mathscr T,M).
\]

11.\quad Sous les conditions du th de Quillen, ($X_0$ de car. nulle, \deleted{\unsure{la}} $(\bigwedge^*\Omega,\delta)$ \deleted{\unsure{loc. libre}}\inserted{avec $\Omega$ plat $\unsure{\Leftarrow|\Longleftrightarrow}$} correspondant à une catégorie formelle $\mathcal C$ avec $\Omega$ plat, pour un Module $M$ sur $\underline O_{X_0}$, les structures de $C$-descente sur $M$ correspondent aux \deleted{\unsure{struct.}} $\Omega$-connexions sur $M$.

% slanted
\margin{Généralise mon th. fondamental de démarrage des cristaux\,\ldots. Peut-il s'y ramener, en raisonnant dans un topos convenable ?!}

De plus, la cohomologie de De Rham de $M$, i.e $\mathbb H^{\unsure{\alpha|*}}(X,\bigwedge^{\unsure{\alpha|2|*}}\Omega\otimes M)$, est can. isom à la coh. de Hochschild (ou Čech-Alexander) de $C$ à coeff. \unsure{dans|\deleted{dans}} $M$, plus précisément il y a un quasi-isom \inserted{canonique} de \inserted{complexes de} faisceaux abéliens (\unsure{dans la catégorie|dans 2 catégories} \unsure{dérivée|dérivées})
\[
\bigwedge^*\Omega\otimes M\simeq C^*(\mathcal C,M)
\]

\uline{Variante à puissances divisées}.

% archive sheet number: 30

